% Recuperación documental para el artículo V; RESULTADO_RECUPERADO.
% Propietario: /Users/ruben/Documents/New project/output/REVISION_EDITORIAL_HMT_MD_20260905/01_FUENTES/INTEGRAL/tree/New project/output/ENSAMBLADO_CAUSAL_RESTAURADO_HMT_MD_20260905/fuente/colaboracion/partes_iv_v/tex/residencias/66_correlacion_cuasilibre_entropias_espectrales.tex
% Líneas de origen: 1--94.
% Sólo se adaptan las etiquetas de referencia y, cuando procede, el nivel de títulos.
% La matriz de recuperación conserva la correspondencia y las dependencias.
% Contenido exclusivo del propietario de entropías tipadas.
\subsection{Corrélation quasi libre et entropies spectrales d’occupation}
\label{v:rec:cuasilibre:section:01}

Pour un état fermionique quasi libre invariant de jauge, conservant le
nombre et défini sur un espace à un corps de dimension finie, soit
\(C_{ij}=\langle a_j^\dagger a_i\rangle\) l’opérateur de corrélation à un
corps. L’hypothèse d’invariance de jauge exclut les corrélations anormales
d’appariement, de sorte que \(C\) détermine l’état quasi libre. Si
\(0<C<I\), son générateur modulaire à un corps et la reconstruction réciproque
du corrélateur sont
\begin{equation}
 \boxed{
 K_1=\log\bigl((I-C)C^{-1}\bigr),
 \qquad
 C=(I+e^{K_1})^{-1}.}
 \label{v:rec:eq:puente-cuasilibre-modular-rev3}
\end{equation}

\begin{corollary}[Forme logistique du spectre modulaire]
\label{v:rec:cor:forma-logistica-modular-rev3}
Si \(K_1u_i=\kappa_i u_i\), l’occupation du mode \(u_i\) est
\[
 \langle n_i\rangle=\frac1{1+e^{\kappa_i}}.
\]
Pour \(\kappa_i=\beta(\varepsilon_i-\mu)\), on retrouve la loi de
Fermi--Dirac.
\end{corollary}

\begin{proof}
La diagonalisation spectrale de \(K_1\) diagonalise toute fonction de cet opérateur.
L’application de \(x\mapsto(1+e^x)^{-1}\) à chaque valeur propre produit la formule.
\end{proof}

Si \(C\) a pour valeurs propres \(p_i\in[0,1]\), l’entropie d’occupation de
l’état fermionique quasi libre est, avec la convention \(0\log0=0\),
\begin{equation}
 S_F=-\sum_i\bigl[p_i\log p_i+(1-p_i)\log(1-p_i)\bigr].
 \label{v:rec:eq:entropia-cuasilibre-fermi-rev3}
\end{equation}
Pour un état bosonique gaussien diagonal avec un nombre fini de modes et
des occupations moyennes \(\overline n_i\),
\begin{equation}
 S_B=\sum_i\bigl[(1+\overline n_i)\log(1+\overline n_i)
                  -\overline n_i\log\overline n_i\bigr].
 \label{v:rec:eq:entropia-gaussiana-bose-rev3}
\end{equation}
La même formule s’étend à une collection dénombrable lorsque la série du
membre de droite converge.
Ces entropies quantifient des configurations d’occupation. L’entropie
d’intrication a un autre domaine : l’état réduit induit par une
décomposition tensorielle. L’identification entre les deux requiert de déclarer la
réduction et l’opérateur de corrélation qui la réalisent.

Enfin, pour un état réduit \(\rho_A\), l’\emph{entropie de von
Neumann}
\[
 S_{\rm vN}(\rho_A)=-\operatorname{Tr}(\rho_A\log\rho_A)
\]
est une propriété spectrale de l’opérateur densité. Lorsque \(\rho_A\)
provient d’une purification globale, elle quantifie l’intrication à travers
la coupure. Dans la construction de bord HMT, l’opérateur d’aire, le
Hamiltonien modulaire et la réduction correspondante sont construits dans la
section~\ref{v:ant:sec:area}. La purification produit de la
proposition~\ref{v:ant:prop:ii-area-purificacion-producto} identifie
la bipartition, et l’identité~\eqref{v:ant:eq:ii-area-ley-finita}
exprime sa relation avec les variations surfaciques.

\Needspace{9\baselineskip}
Les applications entre ces domaines sont exprimées au moyen d’identités explicites.
La décomposition d’une mesure selon les fibres de \(q\) satisfait la règle de
la chaîne
\begin{equation}
 \mathsf H_\mu(\mathbf X)
 =\mathsf H_{q_*\mu}(q(\mathbf X))
  +\mathsf H_{\rm fib}(q;\mu).
 \label{v:rec:eq:cadena-entropia-fibras-rev2}
\end{equation}
\Needspace{5\baselineskip}
Une distribution finie \(p=(p_i)\) se réalise diagonalement comme
\[
 \rho_p=\sum_i p_i|i\rangle\langle i|,
 \qquad S_{\rm vN}(\rho_p)=\mathsf H(p).
\]
Si \(P_{N,E}\) projette sur le sous-espace microcanonique de dimension
\(\Omega(N,E)\), alors
\[
 \rho_{\rm mc}=\frac{P_{N,E}}{\Omega(N,E)},
 \qquad S_{\rm vN}(\rho_{\rm mc})=\log\Omega(N,E)=S_{\rm mc}.
\]
Pour un état pur bipartite,
\[
 \rho_A=\operatorname{Tr}_B
 \bigl(|\Psi_{AB}\rangle\langle\Psi_{AB}|\bigr)
\]
définit la réduction qui mesure l’intrication. Enfin, l’égalité entre
\(h_{\rm top}=\log\rho(A)\) et une entropie métrique est atteinte, pour ce
système primitif, au moyen de la mesure de Parry, qui réalise l’entropie métrique
maximale ; elle n’est pas attribuée à une mesure arbitraire d’histoires. La matrice
\(F_{\rm av}\), son enveloppe à mémoire un et cette mesure sont
construites dans les théorèmes~\ref{thm:incidencia-hojas-anexo}
et~\ref{thm:fibonacci-parry-anexo}, avec leur caractérisation variationnelle
explicite dans la proposition~\ref{v:prop:parry-explicita}.
